\section{Workflow Pengembangan Game di Unity}

\begin{learningoutcomes}
\item Memahami workflow dasar pengembangan game di Unity
\item Mengetahui urutan kerja yang efisien
\item Memahami konsep Scene dan Prefab
\item Mampu mengorganisir project dengan baik
\end{learningoutcomes}

\subsection{Workflow Dasar}

Planning, Asset Creation, Scene Setup, Scripting, Testing (Play mode), Iteration, Build untuk target platform.

\subsection{Konsep Scene}

\begin{definisi}[Scene]
Scene adalah container yang berisi semua GameObject, lighting, dan settings untuk level atau bagian tertentu dari game \cite{UnityManual2024}.
\end{definisi}

Best Practices: Satu scene per level, organize dengan folder dalam Hierarchy, gunakan empty GameObject sebagai organizer, save scene secara teratur.

\subsection{Konsep Prefab}

\begin{definisi}[Prefab]
Prefab adalah template GameObject yang dapat digunakan kembali dan diinstansiasi berkali-kali \cite{UnityManual2024}.
\end{definisi}

Gunakan Prefab untuk: GameObject yang digunakan berulang kali, konfigurasi kompleks, atau yang perlu diupdate secara konsisten.

\subsection{Organisasi Project}

Struktur folder yang disarankan: Scenes/, Scripts/, Art/, Audio/, Prefabs/, Materials/, Animations/.

\subsection{Version Control}

Unity bekerja dengan Git. Gunakan .gitignore untuk Unity. Jangan commit Library/ dan Temp/. Commit Assets/ dan ProjectSettings/.
